\chapter{The Machine Measures Its Own Measurement}

\section{Self-application}

The coupling came back a bracket, read as a self-energy. This chapter dwells on the self in ``self-energy'' --- on the fact
that the instrument, in reading the coupling, turned its measuring on its own act, so that the cost of the reading became the
reading itself. Self-application is not a curiosity attached to the climax; it is the climax's mechanism. The coupling is the
price of the instrument measuring the electron, and the electron the instrument measures includes the instrument's own act of
measuring it. This section makes that circle exact and shows it is not vicious but the very structure the reading required.

The physics where this is precise is renormalization. The electron a bench measures is never the bare electron; it is the
\emph{dressed} one --- the bare charge together with the cloud of its own field, the self-interaction it cannot be separated
from. What the instrument reads, $e_{\text{measured}}$, is $e_{\text{bare}}$ corrected by the electron's coupling to its own
field, and the correction is not removable because the field is not detachable from the charge~\cite{feynman1948}. So the
measured coupling already contains the coupling: the quantity being read is a function of the very interaction that does the
reading. This is self-application in its exact form --- the observable is defined by the process that observes it, the
reading and the read are the same coupling appearing twice.

That the cost of the reading \emph{is} the reading follows directly. To read the electron, the instrument must couple to it;
the coupling has an energy, the self-energy; and that self-energy is the coupling the instrument set out to read. There is no
outside vantage from which to read the coupling without paying it, because the reading is itself an instance of the coupling.
The instrument cannot weigh the electron's self-interaction from a position free of self-interaction; it weighs it from
inside, at the cost of the very interaction being weighed. So the reading returns the cost of reading, and the cost of
reading is the coupling. The circle is exact: what the instrument measures when it measures the coupling is the price of its
own measuring.

This is why the reading could only ever have been a bracket. A quantity read from outside, with no cost paid, could in
principle be pinned to a point. A quantity that is the cost of its own reading cannot: every attempt to sharpen it spends
more of the very thing being sharpened, and the resolution floor of Chapter 6 is exactly where that self-spending can no
longer be resolved. The self-application is the reason the floor exists at this depth and the reason the coupling is an
interval. An instrument measuring something external to itself might hope for a point; an instrument measuring the cost of
its own measurement can only ever reach a bracket, because the act of reaching is itself part of what is read.

In the two verbs, self-application is the funge and the tange sharing a subject. To read the coupling, the instrument funges
the electron with its own field --- bags the bare charge and its dressing as one dressed object --- and the tange that reads
that object is the instrument's own coupling to it, the very interaction inside the bag. What selects and what is selected
are the same coupling. So the instrument measures its own measurement because the measuring is inside the measured, and the
next section names who bears the cost of that circle: the witness, the probe, the detector that pays, in the currency the
whole instrument runs on.

\section{The witness bears the cost}

Self-application means the reading costs, and a cost must be borne by something. This section names the bearer: the witness,
the probe that couples to the electron to read it, pays the price of the reading in the currency the instrument runs on. There
is no free observation. To witness the coupling is to interact with it, and to interact is to be altered --- the witness bears
the back-reaction, and the back-reaction it bears is the measurement. This is the physical content of the circle: measurement
is an exchange, and the witness is the one who is charged.

The physics is measurement back-reaction. A probe that reads a system must couple to it, and the coupling perturbs both:
the system is disturbed by being read, and the probe is altered by the reading. The disturbance is not a defect of a clumsy
apparatus but a floor imposed by the interaction itself --- the complementarity Heisenberg made exact, $\Delta x\,\Delta p
\gtrsim \hbar$, in which the act of resolving one quantity charges the conjugate one~\cite{heisenberg1927}. The witness bears
this charge. To read the electron's position sharply is to pay in its momentum; to read the coupling is to pay in the
self-energy the reading costs. The bearer of the cost is whatever couples in to do the reading, and it pays in the conjugate
currency the interaction fixes.

The currency is the instrument's own, which is why the cost can be borne at all without an external ledger. The witness pays
in $\hbar$, in the quantum of action, in the same pulse the whole construction has run on --- not in a standard kept
elsewhere. This closes a loop the book has drawn since Chapter 6: the resolution floor was set by the frame's own
temperature, the self-energy denominated in the apparatus's own constants, and now the cost of witnessing is charged in the
same intrinsic units. The instrument does not borrow a currency to pay for its readings; it pays in its own, and the floor of
what it can read is set by the smallest coin that currency has. The witness bears the cost in $\hbar$, and $\hbar$ is the
instrument's own.

There is a discipline to mark, because ``the witness bears the cost'' could be read as an excuse and it is the opposite. It
does not mean the instrument's reading is corrupted by an unavoidable disturbance and so cannot be trusted. It means the
disturbance \emph{is} the reading --- that the cost the witness bears is precisely the quantity being measured, so that
paying the cost is measuring, not spoiling the measurement. The self-energy is not noise the witness suffers on the way to
the real value; it is the real value, borne as a cost. The back-reaction is the signal. This is the same inversion the whole
book has run on: the residue is not the error in the reading, it is the reading, and here the witness bearing the
back-reaction is not losing the measurement but performing it.

In the two verbs, the witness is the tange that pays for the funge it performs. To read the coupling the witness funges
itself with the electron --- bags probe and system into one interacting pair --- and the tange it takes, the characteristic
it selects, is charged back to it in the conjugate currency. The selection is not free; the witness is billed the turn, in
$\hbar$, exactly as every seam of Chapter 4 was billed. So the cost of the reading is borne by the one who reads, in the coin
the instrument mints itself. The last section takes the two descriptions this circle has produced --- the bare and the
dressed, the source and the field --- and shows they are two names for one object, the way $1$ and $0.999\ldots$ are two
names for one number.

\section{Two descriptions, one object}

The circle of self-application produced two descriptions of the electron --- the bare charge and the dressed one, the source
and its field --- and weighed them against each other to read the coupling. This closing section draws the lesson those two
descriptions carry: they are two names for one object, and the two names are the same name. The book's emblem for this is the
plainest identity in mathematics, $1 = 0.999\ldots$ --- one number under two descriptions, exactly equal, no residue between
them --- and the section shows that the electron's two accounts stand in exactly this relation.

Take the identity at face value, because its exactness is the whole point. The symbol $0.999\ldots$ and the symbol $1$ are
two descriptions of a single real number; they are not close, not approximately equal, but \emph{identical},
\[
  0.999\ldots \;=\; \sum_{k=1}^{\infty} \frac{9}{10^k} \;=\; 1 ,
\]
the geometric series summing exactly to one. There is no gap. One description approaches from below through an endless tower
of nines --- the Cauchy tower of Chapter 2, finite conditions climbing --- and the other states the limit outright, and the
tower and the limit are the same number. This is two descriptions, one object, in its cleanest form: a value that has more
than one name, and whose names, correctly read, denote the identical thing.

The electron's two accounts stand in this relation, and that is why weighing them was legitimate. The source-account and the
field-account are two descriptions of one electron --- one from below, building the charge up, one from around, the field it
establishes --- and correctly reckoned they are the same object, the dressed electron, under two names. The self-energy the
last chapter read is not a gap between two \emph{different} things; it is the residue of describing one thing two ways and
finding the descriptions do not trivially coincide in the instrument's finite bookkeeping. Like the nines and the one, the
two accounts denote a single electron; unlike the nines and the one, the instrument's finite resolution cannot verify their
identity to a point, and the un-verified remainder is exactly the bracket. Two descriptions, one object --- exactly equal in
principle, bracketed apart only by the floor.

This reframes the coupling one final way before the jar. The coupling is not a magnitude the electron possesses; it is the
residue of the electron having two descriptions the instrument must weigh against each other. Where the two names coincide
exactly --- as $0.999\ldots$ and $1$ do --- there would be no coupling to read; the coupling is the measure of how the
finite instrument's two accounts fail to close, floored at the resolution below which the failure cannot be resolved. So the
coupling is a property of the \emph{description}, not only of the electron --- the price of naming one object two ways with a
finite instrument. That is the deepest reading the physical gauge offers of what the coupling is, and it is marked as a
reading: the identity of the two accounts is exact; the residue the instrument reads between them is its own finite
bookkeeping, bracketed at its own floor.

In the two verbs, ``two descriptions, one object'' is two tanges of a single funge that ought to close and, in a finite
instrument, does not quite. The electron is the funge --- one object, bagged whole. The source-account and the field-account
are two tanges, two selections of how to name the bag, which in principle denote the identical thing, as $1$ and
$0.999\ldots$ do. The coupling is the remainder between the two tanges that the instrument's floor prevents it from resolving
to zero. So the machine measures its own measurement by naming one electron two ways and weighing the names --- and the
un-closable remainder, floored and bracketed, is the coupling. The last chapter takes this bracket, sets it beside what the
instrument can and cannot claim, and draws the jar.
